\documentclass{llncs2e/llncs}
\usepackage{makeidx}  % allows for indexgeneration

\usepackage[pdftex]{graphicx}
\graphicspath{{images/}}

\usepackage{hyperref}
\usepackage{url} \urlstyle{sf}

\usepackage{color}
\def\todo#1{{\color{red}TODO:\quad#1}}
\def\addref#1{{\color{red}$[$#1$]$}}

\usepackage{datetime}
\newdate{startdate}{15}{05}{2017}
\newdate{enddate}{04}{08}{2017}

\title{Automated Test Data Generation for Dynamically Typed Programming Languages}
\titlerunning{Automated Test Data Generation for Dynamically Typed Programming Languages}

\author{Simon Bihel\inst{1, 2}}
\authorrunning{Simon Bihel}
\institute{Universit\'e de Rennes 1, France \and \'Ecole Normale Sup\'erieure de Rennes, France\\
\email{\href{mailto:simon.bihel@ens-rennes.fr}{simon.bihel@ens-rennes.fr}}}

\begin{document}
\pagestyle{headings}  % switches on printing of running heads
\mainmatter%          % start of the contributions

\maketitle

\hrulefill%

\begin{center}
  \textbf{Supervisor:} Dr.\ Shin Yoo\\
  COINSE Lab, KAIST, South-Korea
\end{center}

\hrulefill%

\begin{center}
  \displaydate{startdate} --- \displaydate{enddate}
\end{center}

\vfill

\centerline{%
    \includegraphics[width=0.3\textwidth]{Logo_ENS_Rennes}
    \hspace{0.05\textwidth}
    \raisebox{1\height}{\includegraphics[width=0.3\textwidth]{Universite_Rennes_1_logo}}
    \hspace{0.05\textwidth}
    \includegraphics[width=0.3\textwidth]{KAIST_logo}
    \hspace{0.05\textwidth}
    \includegraphics[width=0.3\textwidth]{RB_NB}
}

\vfill

\begin{abstract}
  \todo{}
Despite popular usage, most of the dynamically typed languages still lack
automated test data generation tools.  Most of the existing tools and
approaches depend critically on static and explicit types, which makes it hard
to port any of them over to dynamically typed languages. In this paper we
present a search-based test data generation tool, NAME, which focuses on
primitive types and lists of Python. NAME not only searches for test input
values but also for appropriate types for each input argument. We evaluate NAME
based on the structural coverage it achieves, the cost of the type search, and
the overhead for the execution of the function under test. We further discuss
the future of automated testing for dynamic and object-oriented languages based
on our experience.
\keywords{software engineering, testing, coverage, experimental}
\end{abstract}
%
\newpage


\section{Introduction}
\todo{}

\addref{}

Testing is an essential part of software development. Due to the lack of formal
specification tests are the pillars to ensure well-behavior of programs, be it
lack of crash or coherent results. Various metrics exist to evaluate how
thoroughly tested a software is. One of these is the \textit{code coverage} that
describes the degree to which the source code of a program is executed when a
particular test suite runs.

Automating the generation of tests is essential to free the programmer from
writing them as is it a time-consuming and error-prone task. Evolutionary
methods like search-based software testing has been of interest for several
decades. In particular, the automated test data generation problem which
consists of generating inputs to explore different kinds of properties for a
given program.

Past works on test generation automation have been focused on statically-typed
languages (more details in Section~\ref{relatedwork}). Having information on the
type of test inputs allows tools to focus on the values of inputs and the depth
of complex structures. Having no knowledge on the inputs induces limitations to
the kind of bugs a tool can detects as well as a heavier cost due to tries with
wrongly typed inputs. Those challenges are described more in depth in
Section~\ref{challenges}.

This paper presents a method to automatically test dynamically-typed languages
(with minimal requirements on the properties of the language)
(Section~\ref{contribution}). We implemented it as a tool call NAME for the
Python language. An evaluation can be found in Section~\ref{evaluation}. Many
problems are still left with for example the generation of objects as dynamic
languages are often Object-Oriented, and future paths of work are put down in
Section~\ref{futureworks}.


\section{Related Works}
\label{relatedwork}
\todo{}~\cite{ducasse2011challenges}~\cite{mairhofer2011search}~\cite{splat}~\cite{splats}.


\section{Challenges}
\label{challenges}
\todo{}


\section{Contribution}
\label{contribution}
\todo{}


\section{Evaluation}
\label{evaluation}
\todo{}


\section{Future Works}
\label{futureworks}
\todo{}


\section{Conclusion}
\label{conclusion}
\todo{}


\section*{Acknowledgment}
Thanks to Shin Yoo and the whole COINSE team that has been very friendly and
welcoming. Thanks also to KAIST for hosting me.

This internship was supported by the R\'egion Bretagne with the \textit{Jeune
\`a l'International} grant as well as the ENS Rennes with the \textit{Aide
Transport pour la Mobilit\'e \`a l'International} grant.

\bibliographystyle{llncs2e/splncs03}
\bibliography{biblio}

\end{document}
